\documentclass[xcolor=dvipsnames]{beamer}
\usepackage[utf8]{inputenc}
\usepackage[T1]{fontenc}
\usepackage{fixltx2e}
\usepackage{graphicx}
\usepackage{longtable}
\usepackage{float}
\usepackage{wrapfig}
\usepackage{soul}
\usepackage{textcomp}
\usepackage{marvosym}
\usepackage{wasysym}
\usepackage{latexsym}
\usepackage{amssymb}
\usepackage{hyperref}
\usepackage{cancel}
\usepackage[percent]{overpic}
\usepackage{listings}
\usepackage{color}
\lstset{ %
  language=Python,
  basicstyle=\ttfamily\tiny,
  emphstyle=\color{red},
  keywordstyle=\color{black}\bfseries,
  identifierstyle=\color{DarkOrchid}\ttfamily,
  commentstyle=\color{Brown}\rmfamily\itshape,
  stringstyle=\color{blue}\slshape,
  showstringspaces=false,
  frame=single,                   % adds a frame around the code
}

\setbeamertemplate{navigation symbols}{}
\useoutertheme{infolines}
\usecolortheme[named=violet]{structure}
\setbeamertemplate{items}[circle]

\DeclareGraphicsExtensions{.pdf,.png,.jpg}
\title[Data Handling]{CCD Acceptance Testing Data Handling}
\author{Brett Viren and Tom Throwe}
\institute[BNL]
{
  Physics Department

  \includegraphics[height=1.5cm]{bnl-logo}
}

\date{\today}

\definecolor{rootpink}{RGB}{255,0,255}
\definecolor{macroyellow}{RGB}{255,215,0}

\hypersetup{
  pdfkeywords={},
  pdfsubject={},
  pdfcreator={pdflatex, beamer and emacs the digital blood, sweat and tears}}

\begin{document}

\maketitle

\begin{frame}
\frametitle{Outline}
\tableofcontents
\end{frame}

\begin{frame}
  \LARGE
  \begin{center}
    Caveat: a work in progress \\
    (but some progress in work).
  \end{center}
\end{frame}

\section{Overview}

\begin{frame}[fragile]
  \frametitle{High Level Data Flow}
  \includegraphics[width=\textwidth]{dataflow}
  \begin{enumerate}
  \item Test station / analyzer produces result files.
  \item All result files are stored ``permanently'' to a well defined disk location.
  \item Specially formatted result summary files are validated, parsed
    and their contents are uploaded to the database.
  \end{enumerate}

  \footnotesize
  By using intermediate files in a common, popular format we can make
  a loose coupling between the test stations and analyzers and the
  database.  This will make supporting the varied test platforms and
  environments simpler.
\end{frame}

\section{Job Harness}

\begin{frame}[fragile]
  \frametitle{Job Harness - Version Control}
  \begin{columns}
    \begin{column}{0.5\paperwidth}
      Assure knowledge of what software was run to produce any given
      result.
      \begin{itemize}
      \item Tests run by a common job harness (Linux and Windows).
      \item Operator enters external parameters defining a run which
        are captured in the output.
      \item Station software from a specific (GIT) commit is checked out.
      \item Test jobs run producing results.
      \item Any post-processing to produce required file format.
      \end{itemize}
    \end{column}
    \begin{column}{0.5\paperwidth}
      \includegraphics[width=\textwidth]{job-harness}
    \end{column}
  \end{columns}

  \vspace{2mm}

  $\rightarrow$Currently at the conceptual level.
\end{frame}

\section{Test Result Files}

\begin{frame}
  \frametitle{Station Results}

  A station / analyzer produces result files in three categories:

  \begin{columns}
    \begin{column}{0.5\paperwidth}
      \footnotesize
      \begin{description}
      \item[Metadata FITS] Describes test software and enumerates results
        files and any auxiliary files.  (well defined schema)
      \item[Result FITS] Result summary files to be parsed and uploaded
        into the DB. (well defined schema)
      \item[Auxiliary] Any additional files to archive. (arbitrary
        formats)
      \end{description}
    \end{column}
    \begin{column}{0.4\paperwidth}
      \includegraphics[width=\textwidth]{files}
    \end{column}
  \end{columns}
  All files are archived, first two types are parsed into LIMS.

\end{frame}

\begin{frame}
  \frametitle{Metadata FITS File} 

  Collects information about all result files.  Here a ``test'' means
  a run on a test station or an offline analysis.  Contains 4 HDUs:

  \begin{enumerate}
  \item Primary HDU, header only.
    \begin{description}
    \item[\texttt{TESTNAME}] a canonical name for the test.
    \item[\texttt{DATE\_OBS}] UTC time stamp for when the test was run.
    \item[\texttt{USERNAME}] Name for the test station operator or analyzer.
    \end{description}
  \item Test's software description consisting of list of main
    programs and their GIT SHA1 commit hash and tag.  Lists are stored
    in a FITS table.
  \item List of FITS files to be parsed by LIMS (see next).  Stored as
    a table of filenames and content SHA1 digest hashes.
  \item List of all auxiliary files to be linked into the LIMS
    database.  Lists stored in same form as \#3.
  \end{enumerate}

\end{frame}


\begin{frame}
  \frametitle{Result FITS File}
  Collects all result data that will be parsed into LIMS's DB.

  \vspace{2mm}

  Follows a well defined schema:
  \begin{enumerate}
  \item Common PrimaryHDU for all tests.
    \begin{description}
    \item[\texttt{EXTNAME}] Canonical name for the test as a whole.
    \item[\texttt{SCHEMAVER}] this HDU's schema version
    \end{description}
  \item Test-specific secondary HDUs with some common cards:
    \begin{description}
    \item[\texttt{EXTNAME}] identify the schema that this HDU follows
    \item[\texttt{SCHEMAVER}] this HDU's schema version
    \end{description}
  \end{enumerate}

  \footnotesize

  Many HDUs for the different types of results will share identical schema elements.

  \vspace{2mm}

  A table/data-only HDU still must provide a standard, minimal set of
  cards in the header.
\end{frame}

\begin{frame}
  \frametitle{Example - JimF's Gain/Noise/ColdSpot}
  This result fits into 5 HDUs:
  \begin{description}
  \item[\texttt{Primary}] Common primary HDU
  \item[\texttt{FileRefs}] Table HDU holding columns for file paths and
    one for file checksum (SHA1) to record input files.
  \item[\texttt{Gains}] Table HDU holding two columns of floats, one
    for each gain measurement type.
  \item[\texttt{Noises}] Table HDU holding two columns of floats, one
    for each noise measurement type.
  \item[\texttt{ColdSpot}] Table HDU holding four columns of ints,
    amp\#, pixel count, spot X and spot Y.
  \item[\texttt{LinRange}] Min/max where each amp is within linear
    specs (to be added)
  \item[\texttt{FullWell}] Point at which each amp ``blows up''. (to
    be added)
  \end{description}

  \vspace{5mm}

  \tiny
  (last two just added yesterday)

\end{frame}


\section{Schema Representation In Python}

\begin{frame}
  \frametitle{Schema Definition in Python}

  Steps to defining and validating a FITS schema for our results.

  \vspace{5mm}

  \begin{enumerate}
  \item Adopt \texttt{pyfits} classes to get valid FITS format by-construction.
  \item Enforce our specific schema by adding a \texttt{validate()} method:
    \begin{itemize}
    \item existing \texttt{pyfits.*HDU} classes (bolt-on) for common
      validation code.
    \item subclass \texttt{pyfits.*HDU} classes for test-specific
      validation code.
    \end{itemize}
  \item Implement a Python module for each test station's results that
    hard-codes its specific schema.
  \end{enumerate}

\end{frame}

\begin{frame}
  \frametitle{Current Status of Implementation}

  Now:
  \begin{itemize}
  \item Some basic validation code exists.
  \item Focused on JimF's example data.
    \begin{itemize}
    \item Schema classes written
    \item Simple validation
    \item Read/write FITS file
    \end{itemize}
  \end{itemize}

  \vspace{2mm}

  Some items still to do on the short term:

  \begin{itemize}
  \item Implement versioning to allow schema evolution, if needed (?)
  \item Emit SQL for table creation.  Ties in to schema evolution.
  \end{itemize}

\end{frame}


\begin{frame}[fragile]
  \frametitle{The LSST CCD Acceptance Testing Results package: \texttt{lcatr}}
  \begin{description}
  \item[\texttt{lcatr.schema}] import all of \texttt{pyfits} and
    bolt-on the \texttt{validate()} methods.
    \begin{itemize}
    \item [$\rightarrow$] code should do ``\texttt{import lcatr.schema}'' before using any \texttt{pyfits} objects.
    \end{itemize}
  \item[\texttt{lcatr.gnc}] definitions for JimF's Gain/Noise/ColdSpots results
  \item[\texttt{lcatr.<test>}] each station will have its own in a similar manner
  \end{description}

  For now, code is maintained at:
  \begin{center}
    \url{https://github.com/brettviren/lcatr}    
  \end{center}
\end{frame}

\begin{frame}[fragile]
  \frametitle{Jim's G/N/C Example in the Schema}

  Construct an unfilled schema (only convention-required content added).

  \begin{lstlisting}[emph={GncInputFilesHDU,GncResult,schema,TableHDU,Column}]
#...
class GncInputFilesHDU(schema.TableHDU):
    def __init__(self, filelist = None, sha1list = None, version = 0):
        super(GncInputFilesHDU,self).__init__()
        self.update_ext_name('FileRefs')
        self.update_ext_version(version)
        fc = schema.Column(name='FileName', format='A64', array=filelist)
        hc = schema.Column(name='SHA1Hash', format='A64', array=sha1list)
        self.columns = schema.ColDefs([fc,hc])
        return
    pass
#...
class GncResult(schema.HDUList):
    def __init__(self):
        super(GncResult,self).__init__(hdus = [
                GncPrimaryHDU(),
                GncInputFilesHDU(),
                GncGainsHDU(),
                GncNoiseHDU(),
                GncColdSpotsHDU(),
                ])
        return
    pass
  \end{lstlisting}
  \footnotesize
  \begin{itemize}
  \item   Hard-code GNC's expectation of what HDUs/Columns/Cards/etc are expected.
  \item   Additional GNC-specific validation code can added to these HDU classes.

  \end{itemize}
\end{frame}

\begin{frame}[fragile]
  \frametitle{Closer look at the \texttt{pyfits} Bolt-ons}

  \begin{lstlisting}[emph={PrimaryHDU_validate,TableHDU_validate,PrimaryHDU,TableHDU,validate}]
# in lcatr.schema
from pyfits import *

def PrimaryHDU_validate(self):
  self.verify()               # pyfits-level verify
  for name in ['EXTNAME','EXTVER']: 
    self.header[name]         # just a lookup to trigger a check
  return
PrimaryHDU.validate = PrimaryHDU_validate  # bolt-on

def TableHDU_validate(self):
  self.verify()
  for name in ['EXTNAME','EXTVER']: 
    self.header[name]
  if not len(self.columns):
    raise ValueError,'TableHDU "%s": no columns' % self.name
  if not all([c.array for c in self.columns]):
    raise ValueError, 'TableHDU "%s": not all columns have arrays' % self.name
  return
TableHDU.validate = TableHDU_validate  # bolt-on
  \end{lstlisting}
  \begin{itemize}
  \item Define \texttt{validate()} method and bolt on to \texttt{pyfits.PrimaryHDU} class.
  \item Require standard name and version cards.
  \item Expect this to expand as we learn more about the data.
  \end{itemize}
\end{frame}


\begin{frame}[fragile]
  \frametitle{Example of Filling the GNC Schema}
  
  Filling from code:
  \begin{lstlisting}
    gnc = GncResult()
    gnc['Gains'].columns[0].array = [2.5122094, 2.8183959, 2.5010307, ...]
    gnc['Gains'].columns[1].array = [2.6737275, 2.7658350, 2.6715354, ...]
    # ...
    gnc.validate()
  \end{lstlisting}

  Some convenience methods can be added for a nicer filling interface,
  but this shows underlying idea.

  \vspace{3mm}

  Filling and validating a file that is already produced:

  \begin{lstlisting}
    import lcatr.schema
    import pyfits
    gnc = pyfits.open("gnc.root")
    gnc.validate()
  \end{lstlisting}
\end{frame}

\section{Questions and Open Issues}

\begin{frame}
  \frametitle{Questions on this concept}
  \begin{itemize}
  \item How many test stations/analyzers work in Python?  
    \begin{itemize}
    \item We will need intermediate scripts to generate compliant
      Results FITS files from the others.
    \end{itemize}
  \item How likely will it be that the schema needs to evolve over
    time?
    \begin{itemize}
    \item If we can guarantee all tests produce the same type of data,
      life is easier.
    \end{itemize}
  \item What else?
  \end{itemize}
\end{frame}

\begin{frame}
  \frametitle{Input Needed}
  \begin{itemize}
  \item Any criticism on this high-level design so far.
  \item A detailed description of the information expected from each
    test station and analyzer.
    \begin{itemize}
    \item To the level of types and quantity of all data.
    \item Developing each result's schema will require iteration.
    \item I will handle producing a unified document covering the
      information produced from all tests based on what I get.
    \end{itemize}
  \item Work with each test operator/analyzer to help them adopt this
    schema for their results.
    \begin{itemize}
    \item I will help integrate the \texttt{lcatr} package into your
      Python-based tests and/or will help write any needed
      generator/converter scripts for tests that do not use Python.
    \end{itemize}
  \item We need a stable, reasonably shallow and well defined disk
    area common to (at least) test stations and database ingest
    client.
  \end{itemize}
\end{frame}

\begin{frame}
  \frametitle{Some Open Issues That Need Thought}

  At the not-yet-even-conceptual level (at least in my own head):

  \begin{itemize}
  \item How to handle test station calibration data.
    \begin{itemize}
    \item Need detailed list of what this data is.
    \item Maybe we can handle in a similar way as the main result
      data.
    \end{itemize}
  \item How to handle environmental (netbotz) information.
    \begin{itemize}
    \item Netbotz appliances apparently have both push (HTTP POST and
      FTP) and pull (web service query) delivery
    \item What type of data must be collected?  What is simply
      archived and what is put in the DB?
    \item Archive frequency?
    \item Assure proper subsequent lookup on timestamp.
      \begin{itemize}
      \item Is NTP running on all test stations and ``netbotz''?
      \item What ongoing assurance that clocks stay synced?
      \end{itemize}
    \end{itemize}
  \item Who writes the collection software and designs the database
    table schema?  
  \end{itemize}
\end{frame}

\end{document}

